\documentclass[]{../../../../tex/classes/styledArticle}
\usepackage{../../../../tex/packages/hebrewSupport}
\usepackage{../../../../tex/packages/mathShortcuts}
\usepackage{../../../../tex/packages/theoremsSupport}

\usepackage{listings}
\lstset{
	language=Python,
	numbers=left,
	breaklines,
	basicstyle=\ttfamily,
	commentstyle=\color{gray},
	prebreak=\textrightarrow,
}

\author{שחר פרץ}
\title{אלגוריתמים $\sim$ \textit{הרצאה 4}}
\begin{document}
	\maketitle
	\subsection*{עצים פורסים מינימליים (עפ''מ)}
	\rmark{באנגלית: \en{Minimum Spanning Trees (MST)}}
	נתון גרף קשיר $G$,לא מכוון, ולכל קשת $e \in E$ משקל $w(e) > 0$. המטרה: לחבר את כל צמתי $G$ ברשת של קשתות עם עלות מינימלית (=סכום המשקלים של הקשתות שבחרנו). 
	\theo{צ''ל ש־$E_0 \subset E$ תת־קבוצה של קשתות כך ש־
	\begin{itemize}
		\item $(V, E_0)$ קשיר
		\item $\sum_{e \in E_0} w(e)$ מינימלי
	\end{itemize}
	הוא עץ. }\begin{proof}
		בגלל שהמשקלים גדולים מאפס, ב־$(V, E_0)$  אין מעגלים, שכן אחרת נוכל להסיר קשת על המעגל ולהקטין את המשקל. 
	\end{proof}
	נרצה למצוא $E_0$ כאלו. משום ש־$(V, E_0)$ הוא עץ, נקרא לבעיה ''חיפוש עץ פורש מינימלי``. נבחין שברגע שמחפשים עץ ניתן לאפשר למשקלים $w(e)$ להיות שליליים או אפס, שכן נוכל להוסיף קבוע $c$ לכל המשקלים, והזזנו את שמקל של כל עץ ב־$c \cdot (\sof{V} - 1)$ כי בכל עץ פורש יש $\sof{V} - 1$ קשתות בדיוק. 
	
	אבחנה: עץ פורש מינימלי איננו בהכרח יחיד. לדוגמה, קליקה כשכל משקליה זהים. 
	
	\theo{אם כל המשקלים שונים זה מזה, יש בגרף עפ''מ יחיד. }
	נוכיח תוך כדי דברים אחרים שנראה היום. 
	
	נראה לפחות שני אלגוריתמים למציאת עפ''מים, שניהם ''בשיטה חמדנית`` (Greedy). למה הכוונה? בכל צעד, נבחר את הקשת הטובה ביותר למצב הנוכחי ונוסיף אותה לעץ עד שנגיע ל־$\sof{V} - 1$ קשתות, ובסוף נטען שהגענו לפתרון אופטימלי. יש דוגמאות שבהן דוגמה חמדנית לא עובדת: 
	
	לדוגמה, bin paceing. הבעיה: להכניס כמות של חבילות לתוך ''ארגזים יותר גדולים`` (ניסוח אחר: מקבלים ר רצף של עבודות בגדלים בין $0$ ל־$1$ ומכניסים אותם למכונות עם גדלים בגודל $1$, כאשר רוצים להתשמש בכמה שפחות מכונות סך־הכל). גישה חמדנית אפשרית תגיד שלהכניס עבודה למגירה הראשונה שעוד יש בה מקום. [יש עוד המון גישות חמדניות, כמו להכניס למכונה עם הכי פחות מקום]. הגישות האלו לא עובדות. [שתי הגישות נשברות בעבור קבלה של חמש פעמים 0.2 ואז חמש פעמים 0.8]. אלגוריתמים חמדנים ישתמשו ב־6 מכונות במקום ב־5. זה לא במקרה – זו בעיה NP קשה (נלמד על זה בסיבוכיות). לא ידוע על אלגוריתם פולינומיאלי שמסוגל לפתור אותה. 
	
	כאשר ניגש לבעיית מציאת עץ פורש מינימלי, ניעזר בשמורה שנגדיר בקרוב. 
	\defi{נקרא לקבוצה $A \subset E$ של קשתות \textit{מבטיחה}, אם קיים עץ פורש מינימלי שמכיל את קשתות $A$. }
	הקבוצה הריקה היא קבוצה מבטיחה, וכמובן שכל עץ פורש מינימלי הוא בעצמו קבוצה מבטיחה. 
	\rmark{הגרף סופי, ולכן יש מספר סופי של עצים פורשים, ולכן יש אכן משקל מינימלי של עץ פורש מינימלי וקיים לפחות אחד כזה בגרף קשיר. }
	
	נרצה להתחיל עם $A = \varnothing$, הקבוצה המבטיחה המינימלית, ובכל שלב להוסיף קשת $e$ כך ש־$A \cup \{e\}$ אכן מבטיחה. 
	
	\defi{בהינתן קבוצה מבטיחה $A$, קשת $e$ תקרא בטוחה אם הוספתה ל־$A$ שומרת על תכונת ההבטחה. כלומר $A \cup \{e\}$ קבוצה מבטיחה. }
	\defi{\textit{חתך} (אנ' cut) הוא חלוקה של $V$ לשתי קבוצות זרות \textbf{ולא ריקות} $(S, V \setminus S)$. }
	\defi{קשת $e$ \textit{חוצה את החתך} אם היא מחברת צומת מ־$S$ אל צומת ב־$V \setminus S$, ונאמר שהיא \textit{חוצה קלה} אם משקלה מינימלי מבין כל הקשתות שחוצות את החתך. }
	נבחין שאם בגרף $G$ קיים חתך עבורו כל הקשתות חוצות [את החתך, לא בין DAG־ים ב־DFS], אזי $G$ דו־צדדי. 
	\theo{נניח $A$ מבטיחה ,ו־$(S, V \setminus S)$ חתך עם התכונה שאף קשת של $A$ אינה חוצה אותו. נניח $e$ חוצה קלה של החתך. אזי $e$ קשת בטוחה. }\begin{proof}
		תהא $A$ מבטיחה, $(S, V \setminus S)$ חתך כנדרש עבור $e$ קשת חוצה קלה. מכיוון ש־$A$ מבטיחה, קיים עפ''מ $T$ שמכיל את קשתות $A$. אם $e \in T$ סיימנו. נניח בשלילה $e \notin T$. נסתכל על $T \uplus \{e\}$. בגרף זה יש $\sof{V}$ קשתות ולכן קיים מעגל ו־$e$ נמצאת על המעגל (כי $T$ הוא עץ). המעגל בהכרח חוצה את החתך פעמיים – פעם אחת דרך $e$, ולפחות פעם נוספת דרך $e'$ כלשהי. משום ש־$e'$ קשת שחוצה את החתך $S$ אז $e' \notin A$ (מהנתון לפיו קשתות $A$ לא חוצות את החתך), ובגלל ש־$e$ חוצה קלה אז $w(e') \ge w(e)$. לכן נוכל להגדיר עץ פורש מינימלי חדש $T' = (T \cup \{e\}) \setminus \{e'\}$, ממשקל לא גבוהה יותר מאשר $T$, ומשקל לא גבוה יותר מאשר $T$ ולכן $e$ בטוחה (אם $w(e') > w(e)$, אז נקבל משקל נמוך יותר בסתירה לכך ש־T עפ''מ). 
	\end{proof}
	נשים לב שבקבוצה מבטיחה $A$ אין מעגלים, כלומר היא יער. בגלל המשפט שהוכחנו לעיל, כל מה שנותר לנו זה למצוא חתך כלשהו בכל צעד. האסטרטגיה של איך לבחור את החתך היא זו שמבדילה בין האלגוריתמים, ומהמשפט לעיל מובטחת נכונות. 
	
	נציג \textbf{אסטרטגיה ראשונה – Prim (1957):} אם $C$ רכיב קשירות של $A$ ו־$e$ חוצה קלה בין $S = V, V \setminus S$ אזי $e$ בטוחה ו־$A \cup \{e\}$ מבטיחה. בריצה ''מגדלים`` רכיב קישות אחד גדול. 
	
	עתה, ל\textbf{אסטרטגיה שנ' Kruskal, 1959): }אם קשת $e$ במשקל מינימלי המחברת שני רכיבי קשירות של $A$, אז $e$ בטוחה, ו־$A \cup \{e\}$ מבטיחה. אם היא מחברת את רכיבי הקשירות $C_1, C_2$ ונבנה חתך $(C_1, V \setminus C_1)$ המקיימת שקשתות $A$ לא חוצות אותו ו־$e$ חוצה קלה. 
	
	
	\subsection*{מימושים}
	ראשית נממש את קרוסקל. 
	נמיין את הקשתות לפי סדר משקל עולה. ננסה להשתמש בהן לפי הסדר. כל קשת $e = (u, v)$ בתורה היא קשת בטוחה אמ''מ היא מחברת שני רכיבי קשירות שונים של $A$ הנוכחית. הפלט של האלגוריתם תלוי בסדר מיון הקשתות (שבירת שוויונות). 
	
	נכונות נובעת ישירות מהנכונות של קשת בטוחה וקב' מבטיחה. 
	
	עולה השאלה, מה הסיבוכיות של האלגוריתם הזה? ראשית כל, צריך למיין בעלות של $\sof{E}\log\sof{E}$. משום ש־$\sof{E} \le \oc(\sof{V}^{2})$ מתקיים מחוקי לוגריתמים $\oc(\sof E \log \sof E) = \oc(\sof E \log \sof V)$. 
	
	שאר הפעולות תלויות במימוש לשנו עבור בדיקה של רכיבי קשירות. נשתמש ב־Union Find: מתחזק אוסף של קבוצות זרות, בהתחלה כל צומת הוא קבוצה משלו. יש לנו ב־Union Find את הפעולות הבאות: 
	\begin{itemize}
		\item הראשון \en{\texttt{make-set($v$)}}: יותר קבוצה חדשה בעבור $\{v\}$. 
		\item השני \en{\texttt{find-set($v$)}}: מה הקבוצה שמכילה את $v$ (מקבל מצביע ל־$v$). 
		\item השלישי \en{\texttt{Union($u, v$)}}: לאחד את הקבוצות שמכילות את $u$ ואת $v$. 
	\end{itemize}
	נבחין שנבצע $\oc(\sof{E})$ פעמים את \texttt{find-set}: עבור כל קשת שתי בדיקות. נבצע בדיוק $\sof V - 1$ פעולות \texttt{Union}. 
	
	מימוש אפשרי: כל קבוצה היא עץ מכוון כלפי השורש (המייצג של הקבוצה) ומחזיקים בשורש את גודל הקבוצה. 
	\begin{itemize}
		\item \texttt{find-set} יקח כעומק $u$ בעץ שלו (מטפסים עד השורש). 
		\item \texttt{union} – מוציאים את השורשים, והופכים אחד מהם להיות הבן של השורש השני. נתלה את העץ עם פחות האיברים (כמו אייכמן נגיד) כבן של העץ עם יותר האיברים (כמו במצריים, שכל הבנים הבכורים היו תלויים). 
	\end{itemize}
	ישנם מימושים יותר נורמליים בפועל שלא ראינו (בגלל אחת מהמלחמות, אבל היינו אמורים להיות) במבני נתונים ומומלץ לעבור עליהם. קל לראות באינדוקציה שמובטח שהסיבוכיות של שתי הפעולות תהיה $\oc(\logn)$ כי עומק עץ המכיל $n$ צמתים מקיים חסום ע''י $\logn$ (שימו לב – אין צורך לבצע רוטציות ובלגנים). למימוש לעיל קוראים Union By Rank/Size. 
	
	אפשר לעשו אופטימיזציה נוספת למבנה הנתונים: בזמן שעולים (ב־\texttt{find}) ניקח את השורש של מי שנתקלנו בו לעבר השורש. התוצאה תהיה עץ יותר שמן, אבל עכשיו פעולות \texttt{find-set} עבור מי שנגרר למעלה יהיו מהירות בהרבה. שתי האופטימיזציות נותנות עלות של $\oc(\sof E \ag(\sof V))$. 
	
	''מעכשיו כשאתה רואה אקרמן הפוך, זה המספר חמש. זה קבוע. כל דבר שאתה תגיד, זה יותר קטן. זה המספר חמש.``
	
	סה''כ $\oc(\sof E \log \sof V)$. 
	
	עתה, ניגש לממש את Prim. מקבל רכיב אחד גדול של $A$ שנקרא לו $C$. בכל פעם מוסיפים לו צומת ע''י חיבור באמצעות קשת קלה ביותר מבין מישהו מ־$C$ למישהו מחוץ ל־$C$. בפועל, האלגוריתם מתחזק מהי הקשת הקלה ביותר שמחברת בינו לבין $C$, כלומר, מיהו השכן הכי ''קרוב`` עליו ב־$C$. נחזיק טור עדיפויות (Priority Queue) עם הקשתות הללו, ובכל רגע נדע להוציא מהטור א תהקשת הקלה ביותר. נסמן ב־$r$ את הצומת הראשון ב־$C$. נחזיק מצביע לכל קודקוד מה הקשת שחיברה אותו לעץ. [תזכורת: בתור עדיפויות יש \texttt{add}, \texttt{extract min} ו־\texttt{decrease key}]. נעשה \texttt{add} בהתחלה עם משקל $\infty$ פרט ל־$v$ עם משקל $0$. אנו עומדים לעשות \texttt{extract min} בהוספת קשת $\oc(\sof{V})$ פעמים ו־\texttt{decrease key} נעשה $\oc(\sof E)$ פעמים בבדיקה של האם קשת משפרת. 
	
	איך בונים את טור העדיפויות? 
	\begin{itemize}
		\item נוכל במקרים רבים להמנע מפיבונאצי ודומיהם, ונשתמש פשוט במערך באורך $\sof V$. זה ייתן סיבוכיות של $\oc(\sof E + \sof V^{2})$. אך במידה והגרף צפוף (מספר הקשתות ריבועי במספר הקודקודים) זה אופטימלי. 
		\item בערימה בינארית נקבל $\oc(\sof E \log \sof V)$. 
		\item בערמת פיבונאצי נקבל $\oc(\sof V \log \sof V+ \sof E)$. 
	\end{itemize}
		
		
	\subsection*{תכונות של עפ''מים}
	\theo{קרוסקל יכול להחזיר כל עץ פורש מינימלי. }\begin{proof}
		הרעיון: במיון של הקשתות ניתן עדיפות בשבירת הקשתות לטובת עפ''מ נתון $T$. ברור שכל עוד בחרנו רק קשתות מ־$T$, אין בעיה. ננניח בשלילה ש־$T$ הוא לא התוצאה. נסמן ב־$e$ את הקשת הראשונה שאיננה מ־$T$ שנבחרה. לפניה נבחרו $e_1 \dots e_i \in T$. נסמן את הקשתות שלא נבחרו מ־$T$ ב־$e_{i + 1} \dots e_{\sof V - 1}$. נבחין ש־$e$ לא סוגרת מעגל עם $\{e_1 \dots e_i\}$, כי אחרת לא היינו בוחרים בה. בנוסף, $w(e) < w(e_j)$ לכל $j > i$ שכן אחרת היינו בוחרים קשת מ־$T$ במקום $e$ (זכרו שהנחנו שאנו נותנים עדיפות לקשתות $T$ במקרה של שוויון). ב־$T \cup \{e\}$ יש מעגל שבו משתתפת $e_j \in T$ עבור $j > i$. נוכל להחליף בין $e$ ו־$e_j$ ולקבל עץ פורש מינימלי במשקל נמוך מ־$T$ הוא $T' = (T \cup \{e\}) \setminus \{e_j\}$ בסתירה להיותו של $T$ עפ''מ. 
	\end{proof}
	
	\theo{תהא $f$ פונקציה עולה ממש מ־$\R_+$ ל־$\R_+$, כלומר $\forall x < y \co f(x) < f(y)$. אזי אם נחליף את משקל של כל קשת $e$ לפי $w$ במשקל $f(w(e))$ נשאר עם אותם עפ''מים מינימליים. כלומר $T$ הוא עפ''מ לפי $w$ אמ''מ $T$ עפ''מ לפי $w'(e) = f(w(e))$. }
	\begin{proof}
		קרוסקל יכול להוציא כל עץ פורש מינימלי, והפעלת $f$ לא משפיע על המיון של הקשתות. 
	\end{proof}
	\theo{לכל זוג עפ''מים יש אותה סדרת משקולות של קשתות. }\begin{proof}
		יהיו עפ''מים $T_1$, $T_2$ עם משקלים $T_1 \co x_1 \le x_2 \le \cdots \le x_{\sof V - 1}$ ו־$T_2 \co y_1 \le \cdots \le y_{\sof V - 1}$. צ''ל $\forall i \in [n] \co x_i = y_i$ נניח בשלילה שקיים $i$ מינימלי עבורו $x_{i + 1} < y_{i + 1}$. נתבונן בפונקציה: 
		\[ f(x) = \begin{cases}
			x & x \le x_{i + 1} \\
			x + 1 & x > x_{i + 1}
		\end{cases} \]
		נבחין ש־$f$ עולה ממש ולכן $w'(e) = f(w(e))$ משמר את הכל. נבחין ש־: 
		\[ w'(T_1) \le x_1 + \cdots + x_{i + 1} + (x_{i + 2} + 1) + (x_{i +3} + 1) + \cdots = w(T_1) + v - i - z \]
		וכן:
		\[ w'(T_2) = u_1 + u_2 + + \cdots  + y_i + (y_{i + 1}) + (y_{i + 2}) = w(T_2) + v - i - 1 \]
		קיבלנו סתירה לכך ש־$w'(T_1) = w'(T_2)$. 
	\end{proof}
	\cola{אם כל המשקלים שונים, יש עפ''מ יחיד. }
	
	
	\subsection*{תרגול}
	\exe{נתון גרף לא מכוון וקשיר. כל קשת צבועה באדום או כחול. רוצים למצוא עפ''מ בעל מספר קשתות מקסימלי של קשתות אדומות. פונקציית המשקל $w \co E \to \R$. }
	
	יש שתי גישות – לשנות את המשקלים ולהפעיל קרוסקל, או במקרה של שוויון לבחור את העץ הפורש עם יותר הקשתות האדומות. 
	\begin{itemize}
		\item פתרון א': במיון הקשתות, עבור קרוסקל, ניתן קדימות בשבירת השוויון לקשתות אדומות. 
		\item פתרון ב': נחסר $\eg$ מהמשקל של הקשתות האדומות, כאשר $\eg$ יבחר כך שגם אם הוא יתווסף לעץ $\sof V - 1$ פעמים, הוא לא יגרום לשני עצים להחליף את הסדר היחסי בינהם. לדוגמה: 
		\[ \eg = \frac{\min\{\sof{w(e_i) - w(e_j)} \mid w(e_i) \neq w(e_j)\}}{\sof V} \]
		ונגדיר את הפונקציה: 
		\[ w'(e) = \begin{cases}
			w(e) & \text{כחולה} \ e \\
			w(e) - \eg & \text{אדומה} \ e
		\end{cases} \]
	\end{itemize}
	נמצא עפ''מ לפי $w'$/לפי המיון החדש וזאת תהיה התשובה. 
	\begin{proof}
		נוכיח שנקבל עפ''מ (לפי המשקל $w$ המקורי) ושהוא בעל מס' מקסימלי של קשתות אדומות. 
		\begin{itemize}
			\item אם נניח בשלילה שקיים עפ''מ $T'$ וכן $w(T') < w(T)$ לפי בחירת $w$ מתקיים $w'(T') < w'(T)$ כי $w'$ לא יחליף סדר משקלים של שני עצים. ולכן $T$ לא עפ''מ לפי $w$ בסתירה. 
			\item נניח בשלילה שקיים עפ''מ $T'$ עם יותר קשתות אדומות. אז $w(T) = w(T')$ כי שניהם עפ''מים. מצד שני, החסרנו משקל $\wg$ כפול מספר הקשתות האדומות ולכן $w'(T') < <w'(T)$ ובסתירה. 
		\end{itemize}
	\end{proof}
	
	\exe{נתון גרף לא מכוון וקשיר. $w \co E \to \R_{+}$ משקלים חיוביים על הקשתות. נרצה למצוא $E' \subset E$ בעלת משקל מינימלי כך ששאר הגרף $G = (V_1, E \setminus E')$ חסר מעגלים. }
	פתרון: $E'$ מינימלית, $E \setminus E'$ חסר מעגלים. נרצה לטעון שהוא עץ (ולא יער). \textbf{טענה: }$E \setminus E'$ הוא עץ. \begin{proof}
		נניח בשלילה ש־$E \setminus E'$ אינו עץ. קיימת קשת שמחברת בין שני עצים בתוך היער, שכן $G$ קשיר. נסתכל על הקבוצה $E' \setminus \{e\}$: מתקיים $w(E' \setminus \{e\}) < w(E')$ וגם $E \setminus (E' \setminus \{e\})$ חסר מעגלים בסתירה לכך ש־$E'$ מינימלית. 
	\end{proof}
	
	לכן, מספיק למצוא עץ פורש מקסימלי $T$ ואז $E' = E \setminus T$. נוכל להשתמש ב־$w'(e) = -w(e)$ ולמצוא עץ פורש מינימלי באמצעות קרוסקל. 
	
	\exe{$G = (V, E)$ לא מכוון, קישר, משקל $w \co E \to \R$ על הקשתות. יהיו $T_1$ עפ''מ ו־$T_2$ עץ פורש \textbf{כלשהו}. המשקלים $T_1 \co a_1 \le a_2 \le \cdots a_{\sof V - 1}$ ו־$T_2 \co b_1 \le \cdots \le b_{\sof V - 1}$. צ.ל. $\forall i \in [\sof V - 1]\co a_i \le b_i$. }
	
	\begin{proof}
		נניח בשלילה שקיים $i$ מינימלי עבורו $a_i > b_i$. נוכל בדומה למקודם להגדיר פונקציה עולה ממש $f$ כך שנקבל סתירה להיות $T_1$ עפ''מ מינימלי בעזרת $w'(e) = f(w(e))$. נגדיר $K = w(T_2) - w(T_1) + 1$. נגדיר: 
		\[ f(x) = \begin{cases}
			x & x \le b_i \\
			x + k & \other 
		\end{cases} \]
		מצפים ש־$w'(T_1) \le w'(T_2)$. אבל: 
		\begin{gather*}
			w'(T_1) \ge w(T_1) + (\sof V - 1 - (i - 1))K = w(T_1) + k(\sof V - i) \\
			w'(T_2) \le w(T_2) + (\sof V - 1 - i)K \\
			\implies w'(T_1) - w'(T_2) \ge w(T_1) + \cancel{K(\sof V - 1)} - w(T_2) + \cancel{(\sof V - i)K} - K = 1
		\end{gather*}\envendproof
	\end{proof}
	
	
	מכיוון שאף אחד לא קרוב ללהיות מפוקס, להלן שתי טענות שתוכלו לקרוא את הוכחתן ברשימות התרגול: 
	\begin{itemize}
		\item \exe{נתון גרף לא מכוון $G = (V, E)$ ו־$w \co E \to \R$, וכן $T_1, T_2$ עפ''מים. נגדיר: 
		\begin{align*}
			E_{1, r} &= \{e \in E(T_1) \mid w(e) \le r\} \\
			E_{2, r} &= \{e \in E(T_2) \mid w(e) \le r\}
			\end{align*}
		אז רכיבי הקשירות של $E_{1, r}$ ו־$E_{2, r}$ זהים לכל $r$. }
		\item \exe{נתון גרף לא מכוון וקשיר $G = (V, E)$ כאשר $w \co E \to \R$ פונקציית משקל, וקבוצת קודקודים $L \subset V$. צ.ל. עץ פורש בעל משקל מינימלי כאשר כל צמתי $L$ הם עלים שלו (לאו דווקא עפ''מ בהגדרתו המקורית). }
	\end{itemize}
	
	חידה: מצאו את העץ הפורש עם הקשת הכי כבדה המינימלית
	
	
	
	
	

	\ndoc
\end{document}